\chapter*{Note on AI Assistance}
\addcontentsline{toc}{chapter}{Note on AI Assistance}

I wrote this book to teach myself the subject. The field moved faster than I
could follow by reading papers one at a time, and the only way I found to hold it
together was to write the account I wanted to read. The chapters exist because I
needed them, and they are ordered the way my own understanding had to be built.

Publishing it is an act of hope rather than authority. If the notes that got me
through the material save someone else the same effort, they have done more than
I originally asked of them. Read it as one practitioner's working understanding,
offered in case it is useful, and not as a settled account from someone who had
finished learning.

That is also why the rest of this note exists. This revision was prepared with
substantial assistance from a large language model, and given that the book argues
for auditable systems, it would be poor form to be vague about how the book itself
was produced.

\vspace{0.5em}
\noindent
\textit{What the model did.}
Working under my direction, the model surveyed the arXiv literature published
between September 2025 and September 2026, screening roughly three thousand
papers against the book's subject matter and recommending a shortlist for
inclusion. It drafted prose throughout, revising the existing text for voice and
removing rhetorical patterns that reviewers had flagged. It expanded three
chapters that existed only as section headings with one-line summaries. It
drafted Chapter~\ref{ch:semantics} in full. It added 138 bibliography entries and
integrated roughly 200 references to work from 2025 and 2026.

\vspace{0.5em}
\noindent
\textit{What was verified, and how.}
Every arXiv identifier added in this revision was checked programmatically
against the arXiv API, confirming that the identifier resolves to the title and
authors cited. That check also surfaced defects in the previous edition. Eighteen
citations referred to real work under invented titles, and one referred to a paper
that appears not to exist. Those have been corrected or replaced, and the
replacements were verified the same way.

Verification of identifiers is not verification of claims. Where the text
attributes a specific finding or number to a paper, that attribution was drawn
from the paper's abstract rather than from a full reading. Readers who intend to
build on a specific result should consult the source.

\vspace{0.5em}
\noindent
\textit{What remains mine.}
The book's thesis, structure, argument, and judgment about what belongs in it are
mine. So is responsibility for every claim it makes, including the ones drafted by
a model and the ones I failed to catch. A tool that drafts text does not transfer
authorship, and it does not transfer accountability.

\vspace{0.5em}
\noindent
\textit{Why this note exists.}
Two reasons. Disclosure of this kind is increasingly expected, and in some venues
required. The better reason is that a book making claims about the reliability of
model-generated work owes its readers an account of its own provenance.

